% GENERATED FILE. Edit canonical lessons, then run npm run generate:books.
% canonical-source-hash: fnv1a64:1488c38a015e054e
% canonical-lessons: ML-C187-visuka, ML-C187-viriyuka, ML-C187-natuka, ML-C187-kuraykkuka, ML-C187-udikkuka

\chapter{To blow, To bloom, To plant, To bark, To rise}
\label{ch:ml-to-blow-to-bloom-to-plant-to-bark-to-rise}
\hlchaptermodality{187}

\begin{chapteropening}
\textbf{By the end of this chapter:} \emph{I can say to blow, to bloom, to plant, to bark and to rise.}

Complete the last lesson of chapter 187: I can say to blow, to bloom, to plant, to bark and to rise.
\end{chapteropening}

\section[\texorpdfstring{\ml{വീശുക} (vīśuka)}{വീശുക (vīśuka)}]{\ml{വീശുക} (vīśuka) — to blow (of wind)}

\label{lesson:ML-C187-visuka}



\begin{quote}
Before the new one: say the Malayalam for to sulk, then the Malayalam for to tease.
\end{quote}

\subsection*{You'll want to know: \ml{വീശുക}}
\textbf{\ml{വീശുക}} — \emph{vīśuka} — \textquotedblleft{}to blow (of wind)\textquotedblright{}.

\begin{grammarlens}[title={What you've built}]
One more word for this chapter.
\end{grammarlens}

\subsection*{Your turn}
\begin{itemize}
\raggedright
  \item \emph{Say it:} \emph{vīśuka}
  \item \emph{Say it:} \emph{vīśuka}, once more
  \item \emph{Say it:} say \emph{vīśuka}
  \item \emph{Recall:} say the Malayalam for to sulk, then the Malayalam for to tease, then say \emph{vīśuka} again
\end{itemize}

\begin{tcolorbox}[breakable,colback=teal!4,colframe=teal!35!black,title={Before you move on}]
What does \ml{വീശുക} mean? (\textquotedblleft{}To blow (of wind)\textquotedblright{}.) Say it once more.
\end{tcolorbox}
\section[\texorpdfstring{\ml{വിരിയുക} (viriyuka)}{വിരിയുക (viriyuka)}]{\ml{വിരിയുക} (viriyuka) — to bloom; to hatch}

\label{lesson:ML-C187-viriyuka}



\begin{quote}
Before the new one: say the Malayalam for to tease, then the Malayalam for to blow.
\end{quote}

\subsection*{You'll want to know: \ml{വിരിയുക}}
\textbf{\ml{വിരിയുക}} — \emph{viriyuka} — \textquotedblleft{}to bloom; to hatch\textquotedblright{}.

\begin{grammarlens}[title={What you've built}]
One more word for this chapter.
\end{grammarlens}

\subsection*{Your turn}
\begin{itemize}
\raggedright
  \item \emph{Say it:} \emph{viriyuka}
  \item \emph{Say it:} \emph{viriyuka}, once more
  \item \emph{Say it:} say \emph{viriyuka}
  \item \emph{Recall:} say the Malayalam for to tease, then the Malayalam for to blow, then say \emph{viriyuka} again
\end{itemize}

\begin{tcolorbox}[breakable,colback=teal!4,colframe=teal!35!black,title={Before you move on}]
What does \ml{വിരിയുക} mean? (\textquotedblleft{}To bloom; to hatch\textquotedblright{}.) Say it once more.
\end{tcolorbox}
\section[\texorpdfstring{\ml{നടുക} (naṭuka)}{നടുക (naṭuka)}]{\ml{നടുക} (naṭuka) — to plant}

\label{lesson:ML-C187-natuka}



\begin{quote}
Before the new one: say the Malayalam for to blow, then the Malayalam for to bloom.
\end{quote}

\subsection*{You'll want to know: \ml{നടുക}}
\textbf{\ml{നടുക}} — \emph{naṭuka} — \textquotedblleft{}to plant\textquotedblright{}.

\begin{grammarlens}[title={What you've built}]
One more word for this chapter.
\end{grammarlens}

\subsection*{Your turn}
\begin{itemize}
\raggedright
  \item \emph{Say it:} \emph{naṭuka}
  \item \emph{Say it:} \emph{naṭuka}, once more
  \item \emph{Say it:} say \emph{naṭuka}
  \item \emph{Recall:} say the Malayalam for to blow, then the Malayalam for to bloom, then say \emph{naṭuka} again
\end{itemize}

\begin{tcolorbox}[breakable,colback=teal!4,colframe=teal!35!black,title={Before you move on}]
What does \ml{നടുക} mean? (\textquotedblleft{}To plant\textquotedblright{}.) Say it once more.
\end{tcolorbox}
\section[\texorpdfstring{\ml{കുരയ്ക്കുക} (kuraykkuka)}{കുരയ്ക്കുക (kuraykkuka)}]{\ml{കുരയ്ക്കുക} (kuraykkuka) — to bark}

\label{lesson:ML-C187-kuraykkuka}



\begin{quote}
Before the new one: say the Malayalam for to bloom, then the Malayalam for to plant.
\end{quote}

\subsection*{You'll want to know: \ml{കുരയ്ക്കുക}}
\textbf{\ml{കുരയ്ക്കുക}} — \emph{kuraykkuka} — \textquotedblleft{}to bark\textquotedblright{}.

\begin{grammarlens}[title={What you've built}]
One more word for this chapter.
\end{grammarlens}

\subsection*{Your turn}
\begin{itemize}
\raggedright
  \item \emph{Say it:} \emph{kuraykkuka}
  \item \emph{Say it:} \emph{kuraykkuka}, once more
  \item \emph{Say it:} say \emph{kuraykkuka}
  \item \emph{Recall:} say the Malayalam for to bloom, then the Malayalam for to plant, then say \emph{kuraykkuka} again
\end{itemize}

\begin{tcolorbox}[breakable,colback=teal!4,colframe=teal!35!black,title={Before you move on}]
What does \ml{കുരയ്ക്കുക} mean? (\textquotedblleft{}To bark\textquotedblright{}.) Say it once more.
\end{tcolorbox}
\section[\texorpdfstring{\ml{ഉദിക്കുക} (udikkuka)}{ഉദിക്കുക (udikkuka)}]{\ml{ഉദിക്കുക} (udikkuka) — to rise (of the sun or moon)}

\label{lesson:ML-C187-udikkuka}



\begin{quote}
Before the new one: say the Malayalam for to plant, then the Malayalam for to bark.
\end{quote}

\subsection*{You'll want to know: \ml{ഉദിക്കുക}}
\textbf{\ml{ഉദിക്കുക}} — \emph{udikkuka} — \textquotedblleft{}to rise (of the sun or moon)\textquotedblright{}.

\begin{grammarlens}[title={What you've built}]
One more word for this chapter.
\end{grammarlens}

\subsection*{Your turn}
\begin{itemize}
\raggedright
  \item \emph{Say it:} \emph{udikkuka}
  \item \emph{Say it:} \emph{udikkuka}, once more
  \item \emph{Say it:} say \emph{udikkuka}
  \item \emph{Recall:} say the Malayalam for to plant, then the Malayalam for to bark, then say \emph{udikkuka} again
\end{itemize}

\begin{tcolorbox}[breakable,colback=teal!4,colframe=teal!35!black,title={Before you move on}]
What does \ml{ഉദിക്കുക} mean? (\textquotedblleft{}To rise (of the sun or moon)\textquotedblright{}.) Say it once more.
\end{tcolorbox}
